\originalsection{附录 B: 流程与数据集细节}\paperlabel{B}{chive-app:pipeline_datasets}
\subsection*{B.1 详细调查流程}\paperlabel{B.1}{chive-app:pipeline_detail}
本附录展开第 \ref{chive-sec:pipeline_description} 节四阶段. 各阶段使用不同模型: 被研究的目标生成行为, 调查智能体筛查、调查和验证, 行为分类器测量发生频率.

\textbf{阶段一: 采样.} 在提示集（附录 \ref{chive-app:prompt_sources}）上运行目标, 评价每提示 30 个回答（便宜的训练运行用 10 个）, 温度 1.0, 使用模型 HuggingFace 默认 top-$p$、top-$k$, 非思考模式输出上限 500 token. 思考运行限制为 1,024 个思考 token 与 2,048 个总生成 token. 阶段三每次重采样回答数固定, 所以基线与反事实批大小相同.

\textbf{阶段二: 筛查.} 智能体读取每提示及全部回答, 按量表给意外程度打 1--5 分; 每次 API 调用还从总计 185 个示例中采五个 few-shot 示例. 调查评分至少 3 的提示. 行为至少发生 30\%, 因低频行为更易受采样噪声影响. 筛查者为每项写简短行为摘要、提交调查智能体的问题（另有第二人称副本供自我解释模型）、以及分类器问题. 调查要测许多重采样中的行为率, 因此此阶段固定一个能由单份回答判断的 yes/no 问题. 此后冻结分类器问题, 全调查不改, 作为测量工具.

筛查阶段在建流程时人工投入最多. 只用量表会产生不有趣、缺少多样性的行为. 我们经过多轮人工查看所标提示, 制作筛查提示和 few-shot 示例.

\textbf{阶段三: 调查.} Opus 智能体拥有对目标的反事实访问. 主工具在编辑提示后重采样, 单轮或多轮皆可, 以阶段一设置返回新回答批. 系统提示有训练与评价两种. 评价的详尽系统提示使智能体运行约 10--15 个实验, 既反驳显著竞争假设, 又确认驱动因素. 较轻训练系统提示用约 5--8 个. 然后必须通过工具调用提交结构化报告, 解释观测行为、核心原因、被驳假设和最终答案. 每个实验还写一行编辑描述及前后实测率.

\textbf{行为分类.} 每次重采样后, 独立分类器 Sonnet 4.6 对批中每份回答回答冻结的问题, 返回一个 yes/no 判决. 行为率是 yes 比例. 同一问题对阶段一回答评分得基线, 每个反事实都报告相对基线的变化. 固定问题评价每批, 所以同次调查实验率可直接比较.

\textbf{阶段四: 验证.} 评判者按 1--10 分评价实验支持解释的程度: 大致 1--2 不支持, 5--6 中等, 7--8 有小问题, 9--10 完全支持. 评价资料平均五个独立评判样本; 训练用一个样本（Qwen3-8B）, 或在单个 vLLM 请求中共享 prefill 抽五个, 成本略高于一个（397B 自我调查, 附录 \ref{chive-app:costs}）. 每个评判检查分类器问题是否公平操作化行为, 描述编辑是否与实际改动一致, 实验是否隔离所称原因. 一个样本可能已足够, 多样本未必值得额外成本. 评价保留评分至少 8 的调查, 训练至少 7.

全部评价数据集的调查者均为 Opus 4.6, 包括 Qwen3.5-397B-A17B 目标. 训练应用另用该模型调查自己, 消除调查者比目标更强的混杂（第 \ref{chive-sec:training_setup} 节）.

\subsection*{B.2 提示来源}\paperlabel{B.2}{chive-app:prompt_sources}
流程接收结束于用户轮次的记录, 不依赖其来源. 两类来源是 WildChat 混合（默认, 全部训练与分布内评价）及 PETRI 审计记录（第 \ref{chive-sec:training_heldout} 节分布外来源）.

\textbf{WildChat 混合.} 默认 WildChat\cite{chive-zhao2024wildchat1mchatgptinteraction} 是真实用户--聊天机器人语料. 再从三项公开集加入携带系统提示或工具使用的智能体记录: Hermes function calling\cite{chive-hermesfunctioncalling2024}、ToolACE\cite{chive-liu2024toolace}、SystemChat-2.0\cite{chive-systemchat2024}. 最终池 91.3\% WildChat, 5.2\% ToolACE, 2.5\% SystemChat-2.0, 1.0\% Hermes. 全文称“WildChat”.

WildChat 有大量几乎相同的自动化提示族, 如共享模板的图像生成、资料标注机器人. 因其开头相同, 按用户文字前 100 字符分组, 忽略空白差异, 同前缀每组至多保留五个. 记录限制为 400--8,000 字符. 太短容易引出简单、泛泛的行为, 反事实调查没什么可解释; 太长删除以降低成本.

\textbf{PETRI.} 分布外来源用 PETRI\cite{chive-petri2025}, 自动审计工具中审计模型引导目标经历与对齐有关的多轮场景, 模拟用户轮次和工具结果. 以 Claude Sonnet 4.6 为审计者, 对每个目标运行 173 个种子场景各两次, 提供合成工具并保留各场景系统提示. 每条记录在每个助手轮次截断, 各截断成为流程的一份提示; 超过 4,096 token 的删除. 得到 Qwen3-8B 的 2,061 份提示、Qwen3.5-397B-A17B 的 1,732 份, 流程不改地运行.

\subsection*{B.3 评价集构造与过滤}\paperlabel{B.3}{chive-app:eval_dataset}
构造三步: 流程产生反事实调查, 每调查抽四个主张, 按质量过滤. WildChat 用固定 4,714 个混合提示, 来自前述池前 5,000 个, 排除也在训练中的 286 个; 各目标相同切片. PETRI 运行则使用同附录的 PETRI 记录.

主张写作两步. 先要求 Opus 4.8 选最能说明原因的两个实验和驳斥合理假设的两个实验. 再由 Sonnet 4.6 描述原记录与反事实记录差异. 它不知道实验结果, 以免措辞泄漏结果; 还标出改变多个变量的混杂反事实.

如第 \ref{chive-sec:pipeline_quality} 节所述, 用三项 LLM 过滤选择原因足够具体、预期可预测的行为. 通过的主张构成正文评价集. 过滤不改结论: 全部受测方法在完整未过滤库上评价, 无增益结果仍在（附录 \ref{chive-app:allruns}）. 发布未过滤全集和代码, 可复现选择、不用过滤、或用更严格过滤.

\textbf{过滤一: 机制具体性.} 很多解释含模糊或模板机制, 如“短故事中 A、B、C 同在时, 模型总选某角色名”. 评判给机制 1 分（模糊统计关联）至 4 分（具体无歧义）, 保留至少 3. 建评价集时应用, 标签还平衡为真/假各 50\%, 至多 2,000 个, 固定每智能体受评主张（表 \ref{chive-tab:eval_dataset_funnel}）.

\textbf{过滤二: 反事实可重复性.} 编辑以意图描述（“移除故事里此角色”“把运动员改为音乐家”）, 而非精确文字 diff, 聚焦语义原因而非精确措辞. 代价是描述可有多种实现, 对微小措辞敏感行为未必重现. 由能运行反事实的智能体重跑, 保留成功重现者, 约留 90\%.

\textbf{过滤三: 单因素干预.} 一些反事实同时改多项, 因调查者合测因素, 或编辑常超过 1,000 token 的长记录时意外改其他东西. 写编辑描述时 LLM 标混杂, 约留 90\%.

过滤集大小随运行变化（表 \ref{chive-tab:eval_dataset_funnel}）. 初始均为 4,714 提示, 最终由通过过滤多少决定. 强模型产生具体原因行为较少, 最终主张也少, Qwen3.5-397B-A17B 最小.

筛查要求行为至少发生 30\% 以限制噪声. 构造正例再要求变化至少 50 个百分点. 因而部分调查无法使用: 消除发生率 30\% 的行为不达阈值. 未来希望覆盖更低频行为, 要保持相同采样误差需每反事实更多回答.

\begin{table}[ht]\centering
\begin{tabular}{lrrrrr}\toprule
 & \multicolumn{3}{c}{WildChat} & \multicolumn{2}{c}{PETRI}\\
阶段 & Gemma27B & Qwen8B & Qwen397B & Qwen8B & Qwen397B\\\midrule
调查 & 1408 & 1320 & 1268 & 625 & 409\\
主张 & 4433 & 4041 & 3972 & 2125 & 1322\\
具体性$\ge3$, 平衡 & 1637 & 1870 & 1309 & 1442 & 708\\
再加可重复 & 1433 & 1666 & 1173 & 1277 & 630\\
再加无混杂 & 1294 & 1497 & 1076 & 1196 & 580\\\bottomrule
\end{tabular}
\caption{评价构造过滤. 每调查至多四个反事实主张, 组成未过滤集. 受评分集应用过滤一（具体性至少3）, 标签各50\%, 至多2000, 每智能体在它上面评分. 正文评价数字再加过滤二（可重复）和三（无混杂）. 两者独立且事后检查, 顺序无影响. 列中模型分别为 Gemma-3-27B、Qwen3-8B、Qwen3.5-397B.}\label{chive-tab:eval_dataset_funnel}
\end{table}

\subsection*{B.4 评价主张格式}\paperlabel{B.4}{chive-app:eval_claim_format}
每个主张写成下面固定自然语言模板. 用 30 个百分点（介于选真/假的 15 与 50 阈值近中点）作客观阈值, 而非模糊的“行为变化”. 两应用格式相同: 可解释性预测器输出 $[0,1]$ 概率并用 AUROC 评分; 自我解释训练目标回答 \texttt{Yes/No}. 字段 \texttt{\{behavior\}} 为分类器问题, \texttt{\{baseline\}} 为实测基线四舍五入的百分比, \texttt{\{intervention\}} 为不知道实测效应的 LLM 写的编辑描述.

\textbf{反事实主张模板（完整翻译）.}
\begin{quote}
行为（补全在满足下述条件时评分为“yes”）: \texttt{\{behavior\}}

基线: 原提示、温度 1.0 下, 此行为约发生于 \texttt{\{baseline\}}\% 的补全.

干预: \texttt{\{intervention\}}

预测: 应用此干预（其他一切不变）, 会使行为发生频率相对约 \texttt{\{baseline\}}\% 的基线改变至少 30 个百分点（任一方向）.

上面的主张为真吗?
\end{quote}
\textbf{反事实主张例子（完整翻译）.} 图 \ref{chive-fig:datapoint} 的 \texttt{randomNum} 范围错误所生成四主张之一:
\begin{quote}
行为（补全在满足下述条件时评分为“yes”）: 回答是否把 \texttt{randomNum} 输出范围描述为约 $[\mathrm{min},\mathrm{max})$ 或“min 与 max 之间”（把 max 参数当输出上界）, 而非正确指出范围为 $[\mathrm{min},\mathrm{max}+\mathrm{min})$ 或等价形式? 若声称范围上限等于或接近 max, 回答 Yes; 若正确指出延伸到 max+min, 或根本没概括范围, 回答 No.

基线: 原提示、温度 1.0 下, 行为约发生于 87\% 的补全.

干预: \texttt{randomNum} 参数名从 max、min 改成 a、b. 其余代码, 包括函数体与全部调用位置, 完全不变.

预测: 应用此干预（其他一切不变）, 会使发生频率相对约 87\% 基线改变至少 30 个百分点（任一方向）.

上面的主张为真吗?
\end{quote}
重跑后行为从 26/30 降到 1/30（83 个百分点变化）, 超过 30 阈值, 因而标签为 Yes.

\subsection*{B.5 数据生成成本}\paperlabel{B.5}{chive-app:costs}
本节成本都对应 Opus 4.6 筛查、调查、验证的运行, 按当时 Anthropic token 价: 每百万输入 5 美元、输出 25 美元、缓存写 6.25 美元、缓存读 0.50 美元. 筛查与验证是单次 API 调用, 用 Batch API 价（50\% 折扣）. 调查是依赖提示缓存的序贯多轮循环, 用标准价. 排除目标服务成本: 阶段一补全与反事实回答都在本地 vLLM 生成, 不经 API.

\begin{table}[ht]\centering
\begin{tabular}{lrr}\toprule
项目 & 评价资料 & 训练资料\\\midrule
每行为率回答数 & 30 & 10\\
阶段四评判样本数 & 5 & 1\\\midrule
筛查（摊销每调查） & \$0.20 & \$0.08\\
调查 & \$0.96 & \$0.70\\
验证 & \$0.75 & \$0.15\\\midrule
每调查合计 & \$1.9 & \$0.9\\
一次运行调查数 & 1736 & 28330\\
整次运行成本 & \$3.3k & \$26k\\\bottomrule
\end{tabular}
\caption{生成资料的每调查 Opus API 成本, 区分评价与训练制度.}\label{chive-tab:datagen_costs}
\end{table}

\textbf{评价资料.} 每次 4,714 提示运行约花 3.3 千美元 API, 得 1,700--2,000 个验证调查, 每个约1.9美元. 调查阶段 token 最多（每调查约470k输入, 每轮重发上下文）, 但缓存显著降低成本: 计费输入约80\%为缓存读, 使每调查成本0.96美元而非无缓存约2.6美元. 验证几乎与调查等贵: 五次独立 Batch 调用各重发约42k token完整验证提示. 如 B.1 所述, 五样本可能不必要, 一个样本可把验证成本减80\%.

\textbf{训练资料.} 更轻系统提示使实验少、行为率用10而非30回答、验证用1而非5样本, 所以便宜. Opus 调查 Qwen3-8B 的训练运行每调查约0.9美元, 28,330个约26千美元（第 \ref{chive-sec:training} 节用约22k调查训练切分）.

\textbf{较便宜调查者.} 不为 Qwen3.5-397B-A17B 自我调查报价, 因自行服务成本强烈取决部署: 自托管还是API、支持提示缓存否、缓存命中率. 各调查者 token 数类似, 因阶段思考预算一致. 使用与 Opus 相同的4,096思考 token预算; 397B调查阶段每调查180k输入、5.1k输出, 相比 Opus训练运行261k计费输入、8.2k输出.
